\documentclass[11pt]{article}
\usepackage[T1]{fontenc}
\usepackage{lmodern}
\usepackage[margin=1in]{geometry}
\usepackage{amsmath,amssymb,amsthm}
\usepackage{microtype}
\usepackage[colorlinks=true,linkcolor=blue,citecolor=blue,urlcolor=blue]{hyperref}
\hypersetup{pdftitle={Failure of exterior observability for a Schrodinger equation with slower surface dispersion},pdfauthor={Kaiwen Jin}}
\newtheorem{theorem}{Theorem}
\newtheorem{lemma}[theorem]{Lemma}
\newcommand{\HH}{\mathcal H}
\newcommand{\EE}{\mathcal E}
\newcommand{\CC}{\mathbb C}
\newcommand{\RR}{\mathbb R}
\newcommand{\norm}[1]{\left\|#1\right\|}

\title{Failure of exterior observability for a Schr\"odinger equation\\
with slower surface dispersion}
\author{Kaiwen Jin\thanks{Email: \href{mailto:yui.ui.siki@gmail.com}{\texttt{yui.ui.siki@gmail.com}}.}}
\date{October 12, 2026}

\begin{document}
\maketitle

\begin{abstract}
We consider a Schr\"odinger field on a planar annulus with a dynamic inner
boundary and a Dirichlet outer boundary. When the surface dispersion
coefficient is one half of the bulk coefficient, observation of the outer
normal derivative fails to control the initial energy at every positive
time. We construct smooth solutions with unit initial energy and
exponentially small observation on each fixed time interval. High angular
frequencies produce modes concentrated near the inner boundary. A
variational construction on a short collar, followed by a comparison in
the graph norm of the evolution operator, turns these localized modes into
exact solutions on the full annulus.
\end{abstract}

\section{The observation problem}

Dynamic boundary conditions allow the trace of a dispersive field to evolve
along the boundary while remaining coupled to the bulk through the normal
derivative. An exterior observability estimate asks whether the initial
energy can be bounded by the flux measured on a separate boundary component.
Such estimates are central to boundary control of the coupled system
\cite{MM}.

Let
\[
\Omega=\{x\in\RR^2:1<|x|<2\},\qquad
\Gamma_1=\{|x|=1\},\qquad \Gamma_0=\{|x|=2\}.
\]
We study complex-valued solutions \(Z=(z,z_\Gamma)\) of
\begin{equation}\label{eq:system}
\begin{cases}
i z_t+\Delta z=0 & \text{in }\Omega,\\
z=0 & \text{on }\Gamma_0,\\
z_\Gamma=z|_{\Gamma_1},\\
i z_{\Gamma,t}+\frac12\Delta_\Gamma z_\Gamma-\partial_\nu z=0
  & \text{on }\Gamma_1.
\end{cases}
\end{equation}
Here \(\nu\) is outward from the annulus, so it points towards the origin
on \(\Gamma_1\), and \(\Delta_\Gamma\) is the Laplace--Beltrami operator on
the inner unit circle. The energy is
\begin{equation}\label{eq:energy}
\norm{Z}_{\EE}^{2}
=\int_\Omega|\nabla z|^2\,dx
+\frac12\int_{\Gamma_1}|\nabla_\Gamma z_\Gamma|^2\,dS.
\end{equation}
The question in \cite{AIM} is whether, for every \(T>0\), a finite constant
\(C_T\) satisfies
\begin{equation}\label{eq:observation}
\norm{Z(0)}_{\EE}^{2}
\le C_T\int_0^T\int_{\Gamma_0}|\partial_\nu z|^2\,dS\,dt
\end{equation}
for every smooth compatible solution.

Mercado and Morales \cite[author preprint, Theorem 1.9 and equation (3.20)]{MM} prove boundary
controllability when the surface coefficient \(\delta\) is larger than
the bulk coefficient \(d\). Their Carleman argument obtains a surface
gradient term with coefficient \(d(\delta-d)\), and does not transfer
unchanged to the present regime \(d=1\), \(\delta=1/2\).
The slower-surface Schr\"odinger regime is identified as an open question
in \cite[Section 5.1]{CMM}, separately from the corresponding negative
result for waves.

\begin{theorem}\label{thm:failure}
There are constants \(c>0\) and \(N\in\mathbb N\), and smooth compatible
solutions \(Z_n\) of \eqref{eq:system}, indexed by integers \(n\ge N\), such
that
\[
\norm{Z_n(0)}_{\EE}^{2}=1,
\qquad
\int_0^T\int_{\Gamma_0}|\partial_\nu z_n|^2\,dS\,dt
\le K_T n^6e^{-2cn}\quad\text{for every }T>0,
\]
where \(K_T<\infty\) is independent of \(n\). Consequently
\eqref{eq:observation} fails for every \(T>0\).
\end{theorem}

For angular frequency \(n\), the surface energy level is \(n^2/2\), whereas
the bulk angular potential near the inner circle is \(n^2/r^2\). The
difference confines a low surface mode to a thin inner layer. We construct
this mode on a collar where the whole bulk potential lies above the surface
level, then cut it off before extending it to the full annulus. A graph
norm estimate controls the outer derivative of the error in the exact
evolution. This permits the construction without isolating an eigenvalue
of the full annulus from its bulk spectrum.

\section{Angular reduction and the outer trace}

Write \(\phi_n(\theta)=(2\pi)^{-1/2}e^{in\theta}\). For \(1<R\le2\), use
the Hilbert space
\[
\HH_R=L^2((1,R),r\,dr)\oplus\CC,\qquad
\norm{(v,\alpha)}_{\HH_R}^{2}=\int_1^R r|v|^2\,dr+|\alpha|^2,
\]
with inner product linear in the second argument. On the form domain
\[
\mathcal V_R=\{(v,v(1)):v\in H^1(1,R),\ v(R)=0\}
\]
define
\begin{equation}\label{eq:form}
q_{n,R}(v)=
\int_1^R\left(r|v'|^2+\frac{n^2}{r}|v|^2\right)dr
+\frac{n^2}{2}|v(1)|^2.
\end{equation}
All endpoint values use the absolutely continuous Sobolev representatives.
The associated sesquilinear form is the polarization of
\eqref{eq:form}. It is densely defined: approximate a bulk function by
smooth functions supported in the interior, and approximate an arbitrary
surface coordinate by a bump at \(1\) with shrinking bulk mass.
For fixed \(n,R\), its form norm is equivalent to the \(H^1\) norm, and
both trace constraints persist under convergence. Thus it is closed and
nonnegative.

The representation theorem for closed forms \cite[Theorem 2.14]{Teschl}
gives a self-adjoint operator \(A_{n,R}\). Its domain and action are
\begin{equation}\label{eq:operator}
\begin{split}
D(A_{n,R})&=\{(v,v(1)):v\in H^2(1,R),\ v(R)=0\},\\
A_{n,R}(v,v(1))&=\left(L_nv,\frac{n^2}{2}v(1)-v'(1)\right),\\
L_nv&=-v''-r^{-1}v'+n^2r^{-2}v.
\end{split}
\end{equation}
Indeed, the weak bulk equation for a represented form functional gives
\((rv')'\in L^2\), hence \(v\in H^2\). Integration by parts then identifies
the surface component:
\[
\int_1^R r\overline{u'}v'\,dr
=-\int_1^R\overline u\,(rv')'\,dr-\overline{u(1)}v'(1),
\qquad u(R)=0.
\]
Conversely this identity shows that each state in the displayed \(H^2\)
domain represents its form functional. There is no additional Robin
condition in this operator domain.

Put \(A_n=A_{n,2}\). Its unitary group \(e^{-itA_n}\) preserves \(D(A_n)\)
and commutes with \(A_n\) \cite[Theorem 5.1]{Teschl}. If
\(U(t)=(v(t),v(t,1))\) solves \(iU_t=A_nU\), then
\(z(t,r,\theta)=v(t,r)\phi_n(\theta)\) solves \eqref{eq:system}. In
particular the boundary equation is
\[
i v_t(t,1)=\frac{n^2}{2}v(t,1)-v_r(t,1),
\]
which has the sign prescribed by \(\partial_\nu=-\partial_r\) at \(1\).
The angular normalization gives the exact identities
\begin{equation}\label{eq:normalization}
\norm{Z}_{\EE}^{2}=q_{n,2}(v),\qquad
\int_{\Gamma_0}|\partial_\nu z|^2\,dS=2|v_r(2)|^2.
\end{equation}
The factor \(2\) comes from \(dS=2\,d\theta\) on the outer circle.

\begin{lemma}\label{lem:trace}
There is a constant \(C\), independent of \(n\ge1\), such that
\[
|v'(2)|\le C\left((1+n^2)\norm{V}_{\HH_2}
+\norm{A_nV}_{\HH_2}\right),\qquad V=(v,v(1))\in D(A_n).
\]
\end{lemma}
\begin{proof}
Positivity of the form and Cauchy--Schwarz give
\[
\norm{v'}_{L^2(1,2)}^2\le q_{n,2}(v)
=\langle V,A_nV\rangle
\le\norm{V}_{\HH_2}\norm{A_nV}_{\HH_2}.
\]
The bulk part of \eqref{eq:operator} implies
\[
\norm{v''}_{L^2(1,2)}
\le C\left(n^2\norm{V}_{\HH_2}
+\norm{v'}_{L^2(1,2)}+\norm{A_nV}_{\HH_2}\right).
\]
For \(g\in H^1(1,2)\), integration yields
\[
g(2)=\int_1^2 g(r)\,dr+\int_1^2(r-1)g'(r)\,dr.
\]
Apply this identity to \(g=v'\), and use
\(\sqrt{ab}\le(a+b)/2\), to obtain the claim.
\end{proof}

\section{A mode confined to an inner collar}

Fix \(b=5/4\), so \(b^{-2}=16/25>1/2\).
The constants below may depend on \(b\) and fixed cutoffs, but are independent
of \(n\). Their values can increase from one estimate to the next.

\begin{lemma}\label{lem:collar}
For each \(n\ge1\), there are \(\lambda_n\in\RR\) and a real
\(u_n\in C^\infty[1,b]\), positive on \([1,b)\), satisfying
\begin{equation}\label{eq:collar}
\begin{gathered}
L_nu_n=\lambda_nu_n,\qquad u_n(b)=0,\qquad
u_n'(1)=\left(\frac{n^2}{2}-\lambda_n\right)u_n(1),\\
\int_1^b r u_n^2\,dr+u_n(1)^2=1,
\end{gathered}
\end{equation}
and
\begin{equation}\label{eq:eigenvalue}
\frac{n^2}{2}\le\lambda_n\le\frac{n^2}{2}+Cn.
\end{equation}
\end{lemma}
\begin{proof}
Minimize \(q_{n,b}\) over the unit sphere in \(\mathcal V_b\), with the
\(\HH_b\) norm. A minimizing sequence has bounded derivative norm, bounded
value at \(1\), and is uniformly bounded and equicontinuous. It has a
uniformly convergent subsequence and a weakly convergent derivative
subsequence. For the latter extraction and lower semicontinuity one may use
\cite[Lemmas 1.12--1.13]{Teschl}. The identity
\(v(r)-v(1)=\int_1^r v'\) identifies the limit derivative. Uniform
convergence retains the trace at \(b\) and the unit mass; weak lower
semicontinuity applies to \(\sqrt r\,v'\). The minimum is therefore attained.

Replacing a minimizer by its modulus does not increase the derivative
energy. One obtains \(|(|v|)'|\le|v'|\) by regularizing the modulus as
\((|v|^2+\varepsilon^2)^{1/2}\) and passing to a weak derivative limit.
We may choose a nonnegative real minimizer \(u_n\). Variation of
\(q_{n,b}-\lambda_n\norm{\cdot}_{\HH_b}^2\) in real and imaginary
directions gives its weak eigenvector equation, hence \eqref{eq:collar}.
The scalar ODE gives smoothness. If \(u_n\) vanished at an interior point,
its derivative there would vanish too; uniqueness for the linear ODE
would force \(u_n=0\). If \(u_n(1)=0\), the boundary equation gives the same
contradiction. Thus it is positive on \([1,b)\).

Because \(r^{-2}\ge b^{-2}>1/2\), the form satisfies
\(q_{n,b}(v)\ge(n^2/2)\norm{(v,v(1))}_{\HH_b}^2\), giving the lower bound
in \eqref{eq:eigenvalue}. For the upper bound take a fixed smooth
\(\eta\), with \(0\le\eta\le1\), equal to \(1\) near \(1\) and zero
near \(b\). The trial function \(\eta(r)e^{-n(r-1)}\) has squared mass at
least \(1\), and its form is at most
\[
\frac{n^2}{2}+bn+\frac{b\norm{\eta'}_\infty^2}{n}+\frac n2.
\]
Here \(\int_1^b e^{-2n(r-1)}dr\le(2n)^{-1}\) bounds both bulk terms.
The variational characterization gives the upper bound.
\end{proof}

\begin{lemma}\label{lem:decay}
For all sufficiently large \(n\),
\[
u_n(1)^2\ge\frac12.
\]
On \(I=[9/8,6/5]\) there is a fixed \(c>0\) such that
\begin{equation}\label{eq:derivatives}
|u_n|\le Ce^{-cn},\qquad
|u_n'|+|u_n''|\le Cn^2e^{-cn},\qquad
|u_n'''|\le Cn^4e^{-cn}.
\end{equation}
\end{lemma}
\begin{proof}
Set \(y_n(r)=r^{1/2}u_n(r)\). The radial equation becomes
\begin{equation}\label{eq:potential}
y_n''=Q_n y_n,\qquad
Q_n(r)=\frac{n^2-1/4}{r^2}-\lambda_n.
\end{equation}
By \eqref{eq:eigenvalue} and \(b^{-2}-1/2=7/50\), for all large \(n\),
\[
Q_n(r)\ge\kappa^2n^2\quad(1\le r\le b),\qquad
\kappa^2=\frac7{100}.
\]
Compare \(y_n\) with
\[
h_n(r)=y_n(1)\frac{\sinh(\kappa n(b-r))}
{\sinh(\kappa n(b-1))}.
\]
The function \(y_n-h_n\) has zero endpoint values and
\[
(-\partial_r^2+\kappa^2n^2)(y_n-h_n)
=(\kappa^2n^2-Q_n)y_n\le0.
\]
A positive interior maximum is impossible. Since \(y_n(1)=u_n(1)\le1\),
we obtain
\begin{equation}\label{eq:decay}
0\le y_n(r)\le
\frac{e^{-\kappa n(r-1)}}{1-e^{-2\kappa n(b-1)}}
\le Ce^{-\kappa n(r-1)}.
\end{equation}
Consequently
\[
\int_1^b r u_n^2\,dr=\int_1^b y_n^2\,dr\le\frac Cn,
\]
and the normalization in \eqref{eq:collar} gives
\(u_n(1)^2\ge1-C/n\ge1/2\).

To estimate derivatives, take \(h=1/40\). For \(r\in I\), the interval
\([r,r+h]\) lies inside \((1,b)\) and remains at distance at least \(1/8\)
from \(1\). On these intervals \(Q_n,Q_n'=O(n^2)\), so
\eqref{eq:potential} and \eqref{eq:decay} bound \(y_n''\). The identity
\[
y_n'(r)=\frac{y_n(r+h)-y_n(r)}h
-\frac1h\int_r^{r+h}(r+h-s)y_n''(s)\,ds
\]
then bounds \(y_n'\) by \(Cn^2e^{-\kappa n/8}\).
Finally \(y_n'''=Q_n'y_n+Q_ny_n'\) is bounded by
\(Cn^4e^{-\kappa n/8}\). Differentiating \(u_n=r^{-1/2}y_n\) gives
\eqref{eq:derivatives} with \(c=\kappa/8\).
\end{proof}

\section{Exact solutions with small exterior observation}

\begin{proof}[Proof of Theorem~\ref{thm:failure}]
Choose a fixed smooth cutoff \(\chi\), with \(0\le\chi\le1\), equal to
\(1\) on \([1,9/8]\) and zero on \([6/5,b]\). Define \(w_n=\chi u_n\)
on \([1,b]\), extend it by zero to \([1,2]\), and set
\[
W_n=(w_n,u_n(1))\in\HH_2.
\]
The inner boundary equation is unchanged by this cutoff. From
\eqref{eq:collar},
\begin{equation}\label{eq:residual}
F_n=(A_n-\lambda_n)W_n=(f_n,0),\qquad
f_n=-2\chi'u_n'-(\chi''+r^{-1}\chi')u_n.
\end{equation}
The function \(f_n\) is smooth and supported in \(I\), away from both
actual boundaries. Lemma~\ref{lem:decay} gives
\begin{equation}\label{eq:residualbounds}
\norm{F_n}_{\HH_2}\le Cn^2e^{-cn},\qquad
\norm{A_nF_n}_{\HH_2}\le Cn^4e^{-cn}.
\end{equation}
For the second bound, two derivatives of \(f_n\) involve profile derivatives
only up to \(u_n'''\). The potential term contributes \(n^2r^{-2}f_n\);
the surface component of \(A_nF_n\) is zero.
The retained surface mass also gives
\begin{equation}\label{eq:initiallower}
q_{n,2}(w_n)\ge\frac{n^2}{2}u_n(1)^2\ge\frac{n^2}{4}.
\end{equation}

Every successive operator image of \(F_n\) remains smooth and compactly
supported in the interior. Thus \(F_n\) belongs to every \(D(A_n^k)\).
Since \(W_n\in D(A_n)\), the identity
\(A_nW_n=\lambda_nW_n+F_n\) proves inductively that \(W_n\) also belongs
to every \(D(A_n^k)\).
Define the exact evolution
\[
U_n(t)=e^{-itA_n}W_n=(v_n(t),v_n(t,1)).
\]
The unitary group preserves all power domains and gives arbitrary time
regularity in their graph norms. The bulk differential equation in
\eqref{eq:operator} successively gives \(H^{2k}\) spatial regularity from
membership in \(D(A_n^k)\). Finite-interval Sobolev representatives then
give a smooth radial function. The power-domain trace constraints supply
all boundary compatibility identities. Multiplying by the normalized angular factor
\(\phi_n(\theta)=(2\pi)^{-1/2}e^{in\theta}\) gives the smooth compatible
solution \(\widetilde z_n(t,r,\theta)=v_n(t,r)\phi_n(\theta)\) of
\eqref{eq:system}, together with its inner trace.

Compare it with the approximate stationary evolution and write
\[
D_n(t)=U_n(t)-e^{-it\lambda_n}W_n.
\]
Differentiating
\(e^{-i(t-s)A_n}e^{-is\lambda_n}W_n\) in the graph norm gives Duhamel's
formula
\[
D_n(t)=-i\int_0^t e^{-i(t-s)A_n}e^{-is\lambda_n}F_n\,ds.
\]
Unitarity and commutation with \(A_n\) yield
\[
\norm{D_n(t)}_{\HH_2}\le t\norm{F_n}_{\HH_2},\qquad
\norm{A_nD_n(t)}_{\HH_2}\le t\norm{A_nF_n}_{\HH_2}.
\]
As \(w_n\) vanishes near \(2\), its outer derivative is zero.
Lemma~\ref{lem:trace} and \eqref{eq:residualbounds} therefore imply
\[
|\partial_r v_n(t,2)|\le Ct n^4e^{-cn}\qquad(t\ge0).
\]
Using the exact normalization \eqref{eq:normalization}, we find
\[
\int_0^T\int_{\Gamma_0}|\partial_\nu\widetilde z_n|^2\,dS\,dt
=2\int_0^T|\partial_r v_n(t,2)|^2dt
\le K_T n^8e^{-2cn}.
\]
Finally divide \(\widetilde z_n\) and its trace by
\(\sqrt{q_{n,2}(w_n)}\), which is nonzero by \eqref{eq:initiallower}.
The resulting solution \(Z_n\) has initial energy one and observation at
most \(K_T n^6e^{-2cn}\). This tends to zero for every fixed \(T>0\),
contradicting \eqref{eq:observation} for any solution-independent \(C_T\).
\end{proof}

\paragraph{AI use.}
This manuscript was developed with OpenAI Codex (gpt-6.1-sol), including
the proof construction, literature checks and writing. The checks performed
during preparation were AI checks; independent mathematical review is
requested.

\begin{thebibliography}{9}
\bibitem{MM}
A. Mercado and R. Morales,
\emph{Exact controllability for a Schr\"odinger equation with dynamic
boundary conditions},
SIAM J. Control Optim. \textbf{61} (2023), no.~6, 3501--3525.
\href{https://doi.org/10.1137/21M1439407}{doi:10.1137/21M1439407}.
\href{https://arxiv.org/abs/2301.02573v1}{Author preprint: arXiv:2301.02573v1}.

\bibitem{CMM}
S. E. Chorfi, L. Maniar and R. Morales,
\emph{Controllability and inverse problems for hyperbolic and dispersive
equations with dynamic boundary conditions},
preprint (2025).
\href{https://arxiv.org/abs/2505.14795v1}{arXiv:2505.14795v1}.

\bibitem{Teschl}
G. Teschl,
\emph{Mathematical Methods in Quantum Mechanics: With Applications to
Schr\"odinger Operators}, second edition,
Graduate Studies in Mathematics \textbf{157},
American Mathematical Society, 2014.
\href{https://www.mat.univie.ac.at/~gerald/ftp/book-schroe/schroe2.pdf}{Author's online edition}.

\bibitem{AIM}
AIM problem collection,
\emph{Boundary observation of a Schr\"odinger field with slower surface
dispersion}, Problem 386 (2026).
\href{https://github.com/MColbrook/AIM/blob/c9929805cece044705abaaebe1cebddac1ba11a6/problems/386-slow-surface-schrodinger-observability.md}{Immutable problem statement}.
\end{thebibliography}
\end{document}
